% Generated draft: requires mathematical and source review.
\documentclass[11pt]{article}
\usepackage[margin=25mm]{geometry}
\usepackage{fontspec}
\setmainfont{Latin Modern Roman}
\usepackage{amsmath,amssymb,mathtools,unicode-math}
\setmathfont{Latin Modern Math}
\usepackage{enumitem,needspace,xurl,hyperref}
\setlength{\parindent}{0pt}
\setlength{\parskip}{5pt}
\setlength{\emergencystretch}{3em}
\newcommand{\N}{\mathbb{N}}
\newcommand{\Z}{\mathbb{Z}}
\newcommand{\Q}{\mathbb{Q}}
\newcommand{\R}{\mathbb{R}}
\newcommand{\C}{\mathbb{C}}
\newcommand{\F}{\mathbb{F}}
\newcommand{\A}{\mathbb{A}}
\newcommand{\PP}{\mathbb{P}}
\newcommand{\E}{\mathbb{E}}
\newcommand{\eps}{\varepsilon}
\newcommand{\dd}{\,\mathrm{d}}
\begin{document}
\section*{UnsolvedMath ID: 10400211}
\textbf{Source problem code:} catalogue entry
\subsection{Definitions and mathematical statement}
A trivial link, or unlink, is a union of circles bounding pairwise disjoint embedded disks in \(S^3\). Let \(U_k\) denote the \(k\)-component unlink in the intended link category. The source's \((2,2)\)-move conjecture asserts that every link can be transformed to an unlink by finitely many such moves, allowing ambient isotopy:
\[
 \forall L\ \exists k,r\ \exists L_0,\ldots,L_r:\quad
 L_0=L,\quad L_r=U_k,
\]
with each successive pair related by ambient isotopy or an allowed \((2,2)\)-move/inverse. No a priori bound on \(r\) is required, and \(k\) is not fixed to the original component count without knowing the move convention.

The precise local tangle replacement, crossing signs and endpoint identifications are not included. The term \((2,2)\)-move must therefore be recovered from the source rather than replaced by arbitrary crossing changes, an unverified rational-tangle convention or a similarly named move. Orientation/framing conventions are also absent. No independent addition or removal of trivial components is allowed unless it is part of the actual move definition.

The summary calls the conjecture not obviously resolved and supplies no exact modern theorem or counterexample. The historical conjecture and qualified open/triage label are retained as source information pending primary-literature review.
\subsection{Short English statement}
Can every link be converted to some unlink by a finite sequence of the specified \((2,2)\)-moves? The exact local move is missing, so it must not be replaced by another unknotting operation.
\subsection{Sources}
\begin{itemize}
\item Local registry: \url{data/unsolvedmath/problems.json}, permanent UnsolvedMath ID \(10400211\), source code catalogue entry; supplied statement, background, and research summary, updated 2026-08-06.
\item Ohtsuki, \emph{Problems on invariants of knots and 3-manifolds} (2002), Conjecture 12.5, attributed to Y. Nakanishi and T. Harikae, PDF page 161: \url{https://msp.org/gtm/2002/04/gtm-2002-04-024s.pdf}. Supplied PDF-font-structure extraction accessed 2026-07-29.
\item Internal reference [220], Conjecture 1.59(6), is cited by the statement; its bibliographic entry and local move diagram are not supplied.
\end{itemize}
\end{document}
